\scanpage{191}
$X$ préschém. connexe \deleted{\unsure{finitude} \unsure{a priori} \unsure{quelconque}}

$1^\circ)$ $X_{\mathrm{réd}}$ de car $p>0$. Alors dans le pro-torseur fondamental
$(T_\alpha)$, il n'y a pas en général de sections pour les $T_\alpha$.

\margin{$\unsure{\ast}$}
On dispose uniquement, \deleted{de façon canonique}

\margin{Coh. $l$-adique}
(i) des $\unsure{\xi}_l^{(\alpha)}\in T_\alpha(\unsure{\mathbb Z|\mathbb Q}_l)$,
$l\neq p$, provenant sans doute
de $\unsure{\xi}_l\in T_\alpha(\widetilde{\mathbb Z}_{l\mathbb Z})$.
\margin{hensélisé}

\margin{\fbox{\begin{minipage}{.29\linewidth}
Coh. $p$-adique\\
Structure supplémentaire :\\
opérations \fbox{$\uuline{F}\text{ et }\uuline{V}$}\\
$1/2$ \uuline{linéaires} \uuline{\unsure{relatifs}}\\
\unsure{Frobenius sur $k$}
\end{minipage}}}

\diagram{\[
\begin{tikzcd}[column sep=small,row sep=large]
\text{(ii) d'un }\unsure{\xi}_p^{(\alpha)}\in T_\alpha( &
\mathbb W_p(k) \arrow[dr,"\text{``DeRham-Hodge''}"'] &
\begin{gathered}
)\quad\text{si }\deleted{\text{$X$ déf. sur}}\inserted{\text{$k$ est un corps parfait sur $X$}}\\
\deleted{\text{le corps parfait $k$}}\\
[\ \unsure{\xi}_p\ \text{provenant sans doute}\\
\text{d'un }\unsure{\xi}_p\in T_{\unsure{p|\alpha}}(\widetilde{\mathbb Z}_{p\mathbb Z})\ ]
\end{gathered}
\\
\text{(iii)}[\ ?\quad\text{\uline{réduit}}\ \unsure{\xi}_p^{(\alpha)\unsure{0|\circ}}\in T_\alpha( &
\deleted{\unsure{\mathbb W}} & k)\ ]
\end{tikzcd}
\]}
\margin{hensélisé strict.}

\margin{Coh. Hodge-deRham. Mais ici, il faut \deleted{sans doute}
travailler dans la catégorie dérivée. Tenir compte aussi de la
structure \uline{filtrée} de Hodge-DeRham}

$\overline{\mathrm M}$ si $X=\operatorname{Spec}\mathbb F_q$, il y a en général
des \deleted{$\unsure{\xi}$} \unsure{obstructions} \unsure{non triviales} dans
des $H^{\unsure{2|1}}(\mathbb Q,G)$ $-\, -\, -$

\margin{$\unsure{\ast}$}
On utilise des ``foncteurs fibres'' qui induisent sur
les motifs effectifs de poids $0$ le foncteur
canonique. Alors les objets de $T_\alpha(\mathbb Z_l)$ sont
sans doute uniques\inserted{à isom près}, (th. de Lang), de $\unsure{\overline{\mathrm m}|\simeq}$
ceux de $T_\alpha(\unsure{\mathbb W}_p(k))$ si $k$ fini (ib), \inserted{plus gén.} pour un
$k$ parfait de dim \unsure{ch} $1$ (?). En général, il
faut s'attendre que $T_\alpha(k)$ \deleted{soit}\inserted{ait} plus d'un
\unsure{élém.} à isom près, mais il est possible
que $T_{\unsure{\infty}}(k)=\varprojlim T_\alpha(k)$ n'a qu'un \unsure{élém.} à
isom près, celui donné par la \unsure{théorie|théor.} de DeRham
Hodge.

\margin{\unsure{relation avec|relation au} \unsure{le problème} \unsure{galoisien} ??
Il doit être \uline{\unsure{épuisant|équivalent}}, cf \uline{Manin}}

\margin{190}
